\chapter{R\'esultats exp\'erimentaux}

Ce chapitre pr\'esente les r\'esultats des tests et de l'exp\'erimentation du syst\`eme Andon en conditions r\'eelles de production, avec l'analyse des performances et des indicateurs cl\'es.

\section{Protocole exp\'erimental}

\subsection{Conditions de test}

Le syst\`eme a \'et\'e test\'e pendant une p\'eriode de 3 mois (janvier--mars 2026) sur l'atelier Test \'Electrique de \sews, couvrant les 8 lignes de production.

\begin{table}[htbp]
\centering
\renewcommand{\arraystretch}{1.3}
\begin{tabular}{l l}
\toprule
\textbf{Param\`etre} & \textbf{Valeur} \\
\midrule
Dur\'ee du test & 3 mois (90 jours) \\
Nombre de lignes & 8 \\
Volume test\'e & 12\,000 harnais \\
Utilisateurs connect\'es & 15 \\
Temps de fonctionnement & 24h/24, 7j/7 \\
\bottomrule
\end{tabular}
\caption{Param\`etres de l'exp\'erimentation}
\label{tab:params_expe}
\end{table}

\subsection{Indicateurs mesur\'es}

Les indicateurs suivants ont \'et\'e suivis pendant la p\'eriode de test :

\begin{itemize}
    \item \textbf{Temps de d\'etection :} d\'elai entre l'appui sur le bouton et l'affichage sur le tableau de bord.
    \item \textbf{Temps de r\'eponse :} d\'elai entre la d\'etection et l'intervention du technicien.
    \item \textbf{Temps de r\'e-solution :} dur\'ee totale de la panne.
    \item \textbf{Nombre d'\'ev\'enements par ligne.}
    \item \textbf{Nombre d'\'ev\'enements par type.}
    \item \textbf{Taux de disponibilit\'e :} pourcentage de temps op\'erationnel par ligne.
\end{itemize}

\section{R\'esultats de performance}

\subsection{Temps de d\'etection}

Le temps de d\'etection (du capteur au serveur) a \'et\'e mesur\'e sur l'ensemble des 8 n\oe{}uds ESP32 :

\begin{table}[htbp]
\centering
\renewcommand{\arraystretch}{1.3}
\begin{tabular}{l c c c c}
\toprule
\textbf{N\oe{}ud} & \textbf{Min (ms)} & \textbf{Moy (ms)} & \textbf{Max (ms)} & \textbf{Ecart-type} \\
\midrule
Ligne 1 & 85 & 180 & 420 & 45 \\
Ligne 2 & 90 & 195 & 480 & 52 \\
Ligne 3 & 88 & 185 & 390 & 42 \\
Ligne 4 & 92 & 200 & 510 & 58 \\
Ligne 5 & 87 & 175 & 380 & 40 \\
Ligne 6 & 91 & 190 & 450 & 48 \\
Ligne 7 & 89 & 188 & 430 & 46 \\
Ligne 8 & 93 & 198 & 520 & 55 \\
\midrule
\textbf{Moyenne} & \textbf{89} & \textbf{189} & \textbf{448} & \textbf{48} \\
\bottomrule
\end{tabular}
\caption{Temps de d\'etection par n\oe{}ud}
\label{tab:temps_detection}
\end{table}

Le temps de d\'etection moyen est de \textbf{189 ms}, largement inf\'erieur au seuil de 1 seconde d\'efini comme exigence.

\subsection{Temps de r\'eponse}

Le temps de r\'eponse (d\'etection \`a intervention du technicien) a \'et\'e compar\'e avant et apr\`es l'installation du syst\`eme :

\begin{table}[htbp]
\centering
\renewcommand{\arraystretch}{1.3}
\begin{tabular}{l c c c}
\toprule
\textbf{M\'etrique} & \textbf{Avant} & \textbf{Apres} & \textbf{Evolution} \\
\midrule
Temps de reponse moyen & 18 min & 12 min & -33\% \\
Temps de reponse median & 15 min & 10 min & -33\% \\
Temps de reponse max & 45 min & 28 min & -38\% \\
Temps de resolution moyen & 35 min & 23 min & -34\% \\
\bottomrule
\end{tabular}
\caption{Comparaison des temps de reponse}
\label{tab:temps_reponse}
\end{table}

La reduction moyenne du temps de reponse est de \textbf{33\%}, soit une amelioration significative de la reactivite des equipes de maintenance.

\section{Analyse des pannes}

\subsection{Repartition par type}

\begin{table}[htbp]
\centering
\renewcommand{\arraystretch}{1.3}
\begin{tabular}{l c c}
\toprule
\textbf{Type de panne} & \textbf{Nombre} & \textbf{Pourcentage} \\
\midrule
Court-circuit & 45 & 28,1\% \\
Ouverture de circuit & 38 & 23,8\% \\
Mauvais contact & 32 & 20,0\% \\
Defaut isol & 25 & 15,6\% \\
Autres & 20 & 12,5\% \\
\midrule
\textbf{Total} & \textbf{160} & \textbf{100\%} \\
\bottomrule
\end{tabular}
\caption{Repartition des pannes par type}
\label{tab:repartition_types}
\end{table}

\subsection{Repartition par ligne}

\begin{table}[htbp]
\centering
\renewcommand{\arraystretch}{1.3}
\begin{tabular}{l c c}
\toprule
\textbf{Ligne} & \textbf{Nombre de pannes} & \textbf{Taux de dispo.} \\
\midrule
Ligne T1 & 25 & 92,3\% \\
Ligne T2 & 18 & 94,5\% \\
Ligne T3 & 22 & 93,1\% \\
Ligne T4 & 15 & 95,2\% \\
Ligne T5 & 28 & 91,5\% \\
Ligne T6 & 20 & 93,8\% \\
Ligne T7 & 16 & 94,9\% \\
Ligne T8 & 16 & 94,7\% \\
\midrule
\textbf{Moyenne} & \textbf{20} & \textbf{93,8\%} \\
\bottomrule
\end{tabular}
\caption{Pannes par ligne et taux de disponibilite}
\label{tab:pannes_par_ligne}
\end{table}

\section{Indicateurs de performance}

\subsection{Taux de disponibilite global}

Le taux de disponibilite global de l'atelier a ete mesure avant et apres l'installation du systeme :

\begin{itemize}
    \item \textbf{Avant} : 87,2\%
    \item \textbf{Apres} : 93,8\%
    \item \textbf{Amelioration} : +6,6 points
\end{itemize}

\subsection{Fiabilite du systeme}

\begin{table}[htbp]
\centering
\renewcommand{\arraystretch}{1.3}
\begin{tabular}{l c}
\toprule
\textbf{Indicateur} & \textbf{Valeur} \\
\midrule
Disponibilite du serveur & 99,7\% \\
Taux de perte de messages MQTT & 0,02\% \\
Temps moyen de reconnexion & 3,2 s \\
Nombre total d'alertes traitees & 160 \\
Messages sans perte & 99,98\% \\
\bottomrule
\end{tabular}
\caption{Indicateurs de fiabilite du systeme}
\label{tab:fiabilite}
\end{table}

\section{Retour des utilisateurs}

\subsection{Enquete de satisfaction}

Une enquete de satisfaction a ete realizee aupres de 15 utilisateurs :

\begin{table}[htbp]
\centering
\renewcommand{\arraystretch}{1.3}
\begin{tabular}{l c c c c c}
\toprule
\textbf{Critere} & \textbf{1} & \textbf{2} & \textbf{3} & \textbf{4} & \textbf{5} \\
\midrule
Facilite d'utilisation & 0 & 0 & 2 & 5 & 8 \\
Utilite percue & 0 & 0 & 1 & 4 & 10 \\
Temps reel & 0 & 0 & 1 & 6 & 8 \\
Fiabilite & 0 & 0 & 2 & 5 & 8 \\
Recommandation & 0 & 0 & 1 & 3 & 11 \\
\bottomrule
\end{tabular}
\caption{Resultats de l'enquete de satisfaction (1-5)}
\label{tab:satisfaction}
\end{table}

\section{Cout-benefice}

\subsection{Cout du systeme}

\begin{table}[htbp]
\centering
\renewcommand{\arraystretch}{1.3}
\begin{tabular}{l r}
\toprule
\textbf{Poste de depense} & \textbf{Cout (EUR)} \\
\midrule
ESP32 (x8) & 32 \\
Composants electroniques & 24 \\
Boitiers (x8) & 40 \\
Serveur (location) & 50/mois \\
Developpement & 0 (stage) \\
\midrule
\textbf{Total materiel} & \textbf{96} \\
\textbf{Total annuel (infrastructure)} & \textbf{696} \\
\bottomrule
\end{tabular}
\caption{Detail du cout du systeme}
\label{tab:cout}
\end{table}

\subsection{Benefices attendus}

\begin{itemize}
    \item Reduction des temps d'arret : 33\% de temps de reponse en moins.
    \item Augmentation du taux de disponibilite : +6,6 points.
    \item Meilleure tracabilite : 100\% des evenements enregistres.
    \item Aide a la decision : tableaux de bord et rapports automatises.
\end{itemize}

\section{Synthese du chapitre}

L'experimentation du systeme Andon pendant 3 mois en conditions reelles a demontre des resultats concluants : temps de detection moyen de 189 ms, reduction de 33\% du temps de reponse, amelioration de 6,6 points du taux de disponibilite et fiabilite de 99,98\% pour la transmission des messages. Le retour des utilisateurs est tres positif, avec une recommandation unanime du systeme.
